% ---------------------------------------------------------------------------------------------------------
% ロボティクスへの含意・限界
% ---------------------------------------------------------------------------------------------------------
\section{ロボティクスへの含意}
\label{sec:robotics}
本研究の対象は人体だが，方法論はそのままロボットに移る．
第一に，\textbf{必要容量マップ}は設計仕様である．3\,cm 縁を両手で捕捉する動作に要る把持力（体重比 2.3--2.7\,BW），
関節トルクの弾性（$\partial\ln J^*/\partial\ln\tau^{\rm cap}\approx$\elastTau），
離手位相の実行可能窓（$\phi_\ell\in[0.25,0.31]$）は，腕渡り・懸垂移動ロボットのアクチュエータとグリッパの
要求仕様をそのまま与える．
第二に，\textbf{双対場}は「どの制約がいつ効くか」を場として保持するので，
同じ物理の異なる個体（質量・リンク長・アクチュエータ容量）への転移が包絡線定理で一次近似でき，
新個体ごとの再学習を減らす．ロボットでは剛体力学が既知なので，双対場 MPC は方策学習なしに直接使える．
第三に，\textbf{価格を報酬にする}整形は，安全制約（把持容量・摩擦錐・可動域）のペナルティ重みを
人手ではなくオラクルの交換比率で決める方法であり，制約付き強化学習のラグランジュ乗数を
物理オラクルから初期化・固定する使い方に相当する．
第四に，環境と最適化器が同一の記号力学を共有する設計（コード生成した縮約物理）は，
学習器を厳密な最適解と同じ土俵で測るための実験基盤として一般に有用である．

\section{限界と今後}
\label{sec:limits}
平面縮約モデルは矢状面のみを扱い，3 次元ひねりモデルによる片手ずつの捕捉負荷の掃引は計算中である．
把持容量は両手合算の絶対値で，指の疲労・皮膚摩擦の個人差を含まない．
双対場学習は現段階で小規模（\dflNok 個の残り最適化，\dflNrows 状態）であり，
閉ループ成功率は反復とデータ量に強く依存する．
costate の学習は価値の学習より本質的に難しく（未知身体で $R^2\approx$\fitPhold），
点ごとの Pontryagin 制御は条件が悪い．短地平線の厳密力学と組み合わせる MPC が必要であった．
今後は，(i)~3 次元モデルの双対場，(ii)~KKT 系の 2 次感度（Riccati 型の価値ヘッセ行列）を用いた
軌道近傍の 2 次双対ラベル，(iii)~実行時に余裕を持つロバスト最適解（離手時刻誤差・順応捕捉のシナリオ）の双対場，
(iv)~MuJoCo 環境への転移を進める．
